\begin{figure*}[t]
\centering
\includegraphics[width=\textwidth]{figures/figure1_paradox.pdf}
\caption{The perturbation--pruning paradox. (a)~Per-message perturbation impact: acknowledgments (ACKs) are the most locally disruptive message type (median similarity $= 0.34$ after removal vs.\ $0.50$ for non-ACK types; $n = 450$ each, Qwen2.5-7B, GSM8K). Lower similarity indicates greater disruption to subsequent dialogue. (b)~Yet bulk removal of \emph{all} acknowledgments preserves task accuracy across configurations. $\bullet$ = single seed; $\blacklozenge$ = multi-seed mean $\pm$ 95\% CI. Green band: $\pm 1$\,pp TOST equivalence; blue band: $-3$\,pp non-inferiority margin. The Llama/StratQA single-seed point ($+5.0$\,pp) reflects high variance under longer reasoning chains ($N{=}238$).}
\label{fig:paradox}
\vspace{-6pt}
\end{figure*}

Large Language Models (LLMs) achieve stronger collective performance when multiple specialized agents exchange natural-language messages to solve a shared task~\citep{du_debate_2024, qian_scaling_2025, zhang_agentprune_2025}. Every such exchange adds a full model call; recent benchmarks report that multi-agent pipelines consume 5--8$\times$ the cost of equivalent single-agent runs~\citep{zhang_agentprune_2025}. Yet no systematic evaluation protocol exists to determine \emph{which} messages are necessary and which can be safely removed. As deployments move toward production, this evaluation gap becomes both a scientific and an engineering imperative. We show that removing any single acknowledgment derails 91\% of subsequent turns, yet removing \emph{every} acknowledgment preserves task accuracy across tested configurations.

Current approaches to communication efficiency span four levels of granularity. Token-level methods compress prompts before transmission~\citep{jiang_llmlingua_2023, chevalier_autocompressors_2023}. Agent-level methods prune entire agents from the collaboration graph~\citep{wang_agentdropout_2025}. System-level methods restructure topologies or schedules~\citep{chen_optima_2025, lin_supervisoragent_2026, li_agp_2025}. A small number of message-level methods filter or compress individual turns~\citep{zeng_s2mad_2025}, but none model the communicative function a message serves or evaluate whether that function predicts safe removal. Theoretical communication bounds for multi-agent systems exist~\citep{rizvi_martel_communication_2026}, yet the gap between abstract capacity results and practical, function-aware evaluation of message necessity remains open.

This gap is exposed by a specific empirical phenomenon (Figure~\ref{fig:paradox}). Acknowledgments, the brief confirmations agents exchange after substantive turns, are the most disruptive single message type under our remove-ablation criterion (\S\ref{sec:perturbation}): remove any one and the dialogue trajectory derails 91\% of the time. Yet bulk removal of \emph{every} acknowledgment leaves task accuracy non-inferior to the baseline ($\delta{=}{-}3$\,pp, TOST $p < 0.001$). We term this disconnect \emph{perturbation-pruning dissociation}. Single-message sensitivity reflects structural coupling between adjacent turns (generation is conditioned on position, not content), not informational necessity. Function-level analysis provides a more direct diagnostic than surface-content heuristics.

We present the first message-level functional analysis protocol for communication in two-agent LLM systems, validated across 18~configurations spanning three model families (7B--72B), 4~reasoning domains, and up to 3~topologies, with 10-seed replication for the main GSM8K configurations. We make three contributions:

\begin{enumerate}[nosep, leftmargin=*]
    \item \textbf{A reusable three-stage protocol} (annotate, perturb, prune; \S\ref{sec:taxonomy_method}--\ref{sec:e2e_pruning}). We release a 21{,}661-message annotated corpus with five-class functional labels across 3~backbones and 4~domains, with human validation ($\kappa = 0.76$) and a zero-shot annotation pipeline requiring no task-specific training data.

    \item \textbf{A statistical hypothesis-testing framework} (\S\ref{sec:statistical_framework}, \ref{sec:impact_gradient}--\ref{sec:e2e_pruning}). We combine TOST equivalence tests, non-inferiority tests, and McNemar's test to separate genuine accuracy changes from noise, to our knowledge, the first systematic application of formal hypothesis tests to multi-agent communication. The framework surfaces \emph{perturbation-pruning dissociation}: function types maximally disruptive under single removal yet harmless under bulk removal constitute safe pruning targets. Acknowledgment-targeted pruning eliminates over a third of inference calls at non-inferior accuracy.

    \item \textbf{Systematic boundary-condition identification} (\S\ref{sec:causal_injection}, \ref{sec:three_modes}). The protocol specifies not only where pruning succeeds but where it fails: MATH-domain tasks, the Qwen3 backbone, and high-temperature Llama configurations resist safe removal. Causal injection experiments reveal backbone-dependent mechanisms: acknowledgments are structurally redundant for Qwen ($p = 0.71$, five seeds) but serve as an implicit error-injection channel for Llama ($\bar{\Delta}{=}{-}3.88$\,pp, $p = 0.003$, five seeds). Mixed-effects models identify three topology-stable communication profiles across backbones. This boundary analysis constitutes a screening tool for the configurations tested: the protocol delineates what must not be pruned.
\end{enumerate}
